\documentclass[11pt]{article}
\usepackage[margin=1.25in]{geometry}
\usepackage{graphicx,booktabs,amsmath,caption,float}
\usepackage[hidelinks]{hyperref}
\captionsetup{font=small}

\title{Decoding and Synchronization of the BIC Test-Beam Runs 3061--3091\\[4pt]\large The four data streams of the August 2026 runs on a common trigger clock}
\author{Tanner Mengel}
\date{October 28, 2026}

\begin{document}
\maketitle

\section*{Introduction}

The BIC test-beam stack is read out by four independent data acquisition streams: two KCU105 boards with
H2GCROC3 chips (KCU0: SFIL1--4, the trigger paddles and the Baby BCAL North end; KCU1: SFIL5), a CAEN DT5742
digitizer (Baby BCAL South and the paddles), and an AstroPix layer. Each stream writes its own RCDAQ file with
its own clock, and no counter is reset in common. Before the streams can be combined, every trigger of every
stream has to be placed on one clock.

This was first done for run 3062. The runs added to disk on 2026-10-27 (beam runs 3061, 3063--3069 and
cosmic runs 3088, 3090, 3091) are processed here with the same method. Every constant is re-measured on each
run. This note describes
\begin{itemize}
  \item the decoding of the raw files into per-trigger ROOT trees, and its checks;
  \item the synchronization of KCU1, CAEN and AstroPix to KCU0, and its checks;
  \item the resulting common trigger table of each run.
\end{itemize}
Run 3062 is processed again with the new run-generic code, as a check against its published constants.

%%RUNS%%

\section*{Decoding}

\subsection*{Raw data}

Every run is split into four streams, \texttt{\{beam,cosmics\}\_\{ROC\_0,ROC\_1,CAEN,AstroPix\}-0000RRRR-SSSS.evt}.
The files roll over at 3\,GB, so long runs have several segments (run 3064: 140 KCU0, 43 KCU1 and 42 CAEN
segments). Each RCDAQ event carries one data packet and a host-clock packet (RTCLOCK, 1\,ns, the time the
buffer was closed):

\begin{center}
\begin{tabular}{llll}
\toprule
stream & packet & one RCDAQ event holds & clock used for synchronization \\
\midrule
KCU0 & 12001 & 200 UDP frames, $\sim$9 triggers & FPGA timestamp, 6.25\,ns, 64 bit \\
KCU1 & 12002 & same & FPGA timestamp, 6.25\,ns, 64 bit \\
CAEN & 2001 & one trigger & group trigger time, 8.533\,ns, 30+30 bit \\
AstroPix & 13001 & a chunk of the SPI stream & on-chip \texttt{ts} (10\,MHz, 8 bit) and \texttt{fpga\_ts} (40\,MHz, 32 bit) \\
\bottomrule
\end{tabular}
\end{center}

\subsection*{KCU streams}

A KCU trigger is 20 samples of 288 channels, i.e.\ 160 lines of 36 channels (8 lines per sample: 4 ASICs
$\times$ 2 halves). The lines arrive in UDP frames, and a trigger can be spread over two RCDAQ events. The
converter (\texttt{kcu\_to\_root}) uses the local fork of the online decoder pool:
\begin{itemize}
  \item samples are assembled by timestamp: a line joins the open sample within 5 ticks, and samples close
        4 sample periods behind the newest timestamp, so lines arriving out of order do not split a waveform;
  \item the bytes of every line are then looked up by (timestamp, link) and decoded field by field, so the
        trees hold the raw words: ADC, TOA, the full 10-bit TOT code, the TOT flags, the frame header, the
        common mode, the calibration channel and the CRC-32 of every line;
  \item every 7th channel of every sample is compared with the decoder's own ADC value.
\end{itemize}
Each trigger becomes one entry of the \texttt{header} tree and of one tree per subsystem
(\texttt{sfil1}--\texttt{sfil4}, \texttt{paddles}, \texttt{bbcal}, \texttt{unused} on KCU0;
\texttt{sfil5}, \texttt{sfil5north}, \texttt{unused} on KCU1). A line lost in the network leaves $-1$ in its 36
channels for that sample.

\paragraph{Channel map.} The map follows the run-log hookup tables. Runs up to 3069 use the 13 August table,
as for run 3062. After 14 August 21:55 CT (runs from 3085 on) the four SFIL5 North summed boards are on KCU0
ASIC2, channels 153, 156, 163 and 167; for these runs they are written to a new tree \texttt{sfil5n}
(cells 1--4). For runs before 3085 the output is bit-identical to the previous converter.

\paragraph{Segment boundaries.} One condor job converts one segment. It also reads the first 6 packets of
the next segment, so the trigger crossing the boundary is written complete. Triggers with a line in the
segment's first packet go to a separate head file; the merge step drops from it the triggers the previous
segment already holds, and chains head and main files in time order.

\paragraph{Stale clock at run start.} The first few triggers of each KCU file carry the FPGA counter from
before the reset at run start; the counter then restarts near zero. Each file therefore has exactly one
backward timestamp step, flagged in \texttt{header.backward}. The synchronization drops the entries before it.

\subsection*{CAEN}

One RCDAQ event is one trigger: 16 channels (Baby BCAL South ch 0--11, paddles P3--P6 ch 12--15) of 256
DRS4 cells at 1\,GS/s. The converter (\texttt{caen\_to\_root}) checks the event size and marker, and compares a
sample of every 50th trigger and the group trigger time with the stock decoder. The two 30-bit group trigger
times are joined into one counter, \texttt{gtt}, that does not wrap over a run. The 32-bit event time tag
steps back by $2^{31}$ and is not used.

\subsection*{AstroPix}

The stream is a sequence of SPI records rather than triggers. The converter (\texttt{astro\_to\_root}) writes
one entry per 11-byte hit record (\texttt{hits}) and keeps every other record (SPI filler and records cut at
a packet end) in \texttt{records}. A hit is valid if it is on layer 1, chip 0--8. The hit count is compared
with the stock decoder.

\subsection*{Processing}

The 269 segment files were converted by condor jobs, one per segment and stream, and merged per stream and
run. Run 3064 is merged in parts of 10 segments: a single file would be $\sim$350\,GB for KCU0, beyond ROOT's
100\,GB tree limit. Each part starts on a de-duplicated head file, so the parts read in order give every
trigger once.

%%DECODE%%

\section*{Synchronization}

All streams are put on the KCU0 clock (6.25\,ns ticks). Every step is checked against a null control: the
stream shifted by 10\,ms, a neighbouring trigger, or shuffled timestamps. The code is
\texttt{sync/run\_sync.py}; it reads only the merged ROOT files.

\paragraph{Triggers.} A KCU trigger is a \texttt{header} entry after the stale-clock prefix. An entry with
fewer than 20 samples lying $k \times 164$ ticks ($k = 1..19$, $\pm 4$) after another entry is a fragment of
that trigger, not a new trigger.

\paragraph{KCU1 to KCU0.} Both boards count the same 160\,MHz clock from different resets. A coarse offset
comes from the host-clock epochs, $L = t_\mathrm{host} - 6.25\,\mathrm{ns} \times \mathrm{ts}$, per board.
The exact offset is the most common value of $\mathrm{ts}(\mathrm{KCU0}) - \mathrm{ts}(\mathrm{KCU1})$ over
all pairs within $\pm 20$\,ms of the coarse value, for 100\,k sampled KCU1 triggers. It is measured again in
ten time slices of the run.

\paragraph{CAEN to KCU0.} The CAEN time is $t_C = 8.533\,\mathrm{ns} \times \mathtt{gtt}$. The coarse offset
uses the 1st percentile of $t_\mathrm{host} - t$ for both streams, because the CAEN buffers reach the host up
to 2\,s late. The run is cut into slices of at most 300 CAEN triggers and at most 5\,s. In the first slice the
coincidence peak of all KCU0--CAEN time differences is searched over $\pm 3$\,s; each later slice is searched
within $\pm 2$\,ms of the extrapolated offset, in 10\,$\mu$s bins. A straight line through the slice offsets
gives the clock drift. A running median of the nearest-neighbour residuals (201 pairs) then removes the slow
wander of a few $\mu$s. The 5\,s limit on a slice matters at low trigger rates: with a drift of 5\,ppm, a
300-trigger slice of a 20\,Hz run would otherwise smear the peak over 75\,$\mu$s.

\paragraph{Common trigger table.} The KCU0 triggers, the KCU1 triggers moved by the offset and the CAEN
triggers on the KCU0 clock are grouped within 16 ticks (100\,ns). Each row of the table is one trigger,
with the entry number in each stream's merged file ($-1$ where the stream missed it).

\paragraph{Paddle cross-check.} The paddles are read by both KCU0 and CAEN, which gives an independent check
of the matching: on every trigger seen by both, the pattern of fired paddles is compared. A paddle fired on
KCU0 if the pedestal minus the minimum of samples 6--19 exceeds 65 ADC, and on CAEN if the pedestal minus the
minimum of samples 20--229 exceeds 200 ADC. The decoy is the CAEN pattern of the neighbouring trigger.

\paragraph{AstroPix to KCU0.} The AstroPix \texttt{fpga\_ts} stamps the SPI readout, not the hit, so it places
a hit only to within the readout latency (tens of ms). The trigger-level handle is the on-chip
\texttt{ts}: in run 3062 it is a 10\,MHz counter locked to the KCU clock,
$\mathtt{ts} = ((K + s)\,//\,16 - \phi) \bmod 256$ for the partner trigger at KCU0 tick $K$, with a fixed
phase $\phi$. Hits are grouped into clusters (\texttt{ts} within $\pm 2$, read within 1\,ms). Then:
\begin{enumerate}
  \item a coarse offset from the readout envelopes (the 1st percentile of $t_\mathrm{host} - t$ per 10\,s, for
        the last record of each RCDAQ event);
  \item a wide scan: for a subset of clusters, every KCU0 trigger within $\pm 60$\,s is histogrammed in time
        offset (10\,ms bins) and phase; true pairs pile up in one cell, random pairs are flat in phase;
  \item a phase scan within $\pm 250$\,ms of that cell for all clusters, against the same scan with
        \texttt{ts} shuffled between clusters;
  \item the readout latency: the time difference of phase-matched pairs minus that of pairs in the
        opposite phase (control).
\end{enumerate}
Step 2 replaces the beam on/off correlation used for run 3062, which fails in the low-rate runs (no visible
beam structure).

%%SYNC%%

\section*{Summary}

%%SUMMARY%%

\end{document}
